\chapter{Polymyalgia Rheumatica}
\index{Polymyalgia Rheumatica}
\label{sec:Polymyalgia Rheumatica}

\section{Overview}
Polymyalgia rheumatica is an inflammatory disorder of the elderly, characterized by pain and stiffness of the shoulder and pelvic girdles, affecting individuals over 50 years of age (peak incidence 70--80 years), with a female predominance of 2--3:1.
\section{Causes}
\section{Risk Factors}

\begin{itemize}[noitemsep]
  \item Genetic predisposition and family history
  \item Advancing age
  \item Lifestyle factors (diet, exercise, smoking, alcohol)
  \item Environmental exposures
  \item Underlying medical conditions
  \item Medication use
\end{itemize}
\section{Signs and Symptoms}

\subsection{Common symptoms}
\begin{itemize}[noitemsep]
  \item You may experience acute or subacute onset of severe morning stiffness lasting over 30 minutes in the shoulders, neck, and hips, with associated muscle pain, constitutional symptoms (low-grade fever, fatigue, weight loss), and difficulty raising arms or rising from a chair, with peripheral arthritis (knees, wrists) occurring in 25--50\%. PMR overlaps significantly with giant cell arteritis (GCA): 15--30\% of PMR patients have concurrent GCA, and 40--60\% of GCA patients have PMR symptoms. 
  \end{itemize}

\subsection{Emergency warning signs}
\begin{itemize}[noitemsep]
  \item Severe or worsening symptoms of polymyalgia rheumatica require immediate medical evaluation
\end{itemize}
\section{Progression and Stages}
The cause is unknown, but genetic associations with HLA-DR4 and environmental triggers have been implicated, with Disease mechanism involving local inflammation of proximal joints and extra-articular synovial structures, infiltration of macrophages and CD4+ T cells, elevated IL-6 levels, and increased acute-phase reactants. The condition may progress through different stages of severity over time. Early stages often involve mild or subtle symptoms that may be overlooked. Moderate stages involve more noticeable symptoms affecting daily function. Advanced stages may involve significant impairment and complications.
\section{Complications}
\begin{itemize}[noitemsep]
  \item Disease progression leading to worsening symptoms
  \item Secondary infections or injuries
  \item Side effects of treatment
  \item Reduced quality of life
  \item Emotional and psychological impact
\end{itemize}
\section{Diagnosis}
Doctors usually diagnose this based on your symptoms and a physical exam, supported by elevated ESR (often $>$50--100 mm/hr) and CRP, with imaging (ultrasound, MRI) showing synovitis, bursitis, and tenosynovitis. Comprehensive medical history and physical examination Laboratory tests to assess organ function and rule out other conditions Imaging studies to visualize affected structures Specialized testing depending on the affected system Consultation with relevant specialists Monitoring of disease progression over time

\begin{itemize}[noitemsep]
  \item Comprehensive medical history and physical examination
  \item Laboratory tests to assess organ function
  \item Imaging studies to visualize affected structures
  \item Specialized testing depending on the affected system
  \item Consultation with relevant specialists
\end{itemize}
\section{Prevention}
\begin{itemize}[noitemsep]
  \item Maintain a healthy lifestyle with balanced diet and exercise
  \item Avoid known risk factors
  \item Attend regular medical check-ups for early detection
  \item Follow recommended screening guidelines
\end{itemize}
\section{Treatment Options}
Treatment typically involves with low-to-moderate dose corticosteroids (prednisone 15--25 mg/day), typically with rapid clinical response within 24--48 hours and slow tapering over 12--18 months. Medications to manage symptoms and address underlying mechanisms Lifestyle modifications and self-care measures Physical therapy and rehabilitation Surgical intervention when indicated Regular monitoring and follow-up care Patient education and support services

\begin{itemize}[noitemsep]
  \item Medications to manage symptoms and address underlying mechanisms
  \item Lifestyle modifications and self-care measures
  \item Physical therapy and rehabilitation
  \item Surgical intervention when indicated
  \item Regular monitoring and follow-up care
\end{itemize}
\section{Living with It}
\begin{itemize}[noitemsep]
  \item Follow your treatment plan as prescribed
  \item Maintain regular follow-up with your healthcare provider
  \item Make necessary lifestyle adjustments
  \item Seek support from family, friends, or support groups
  \item Stay informed about your condition
\end{itemize}
\section{Outlook}
The outlook is good, with most patients achieving remission within 1--3 years. The outlook varies depending on the specific condition, severity at diagnosis, and response to treatment. Early intervention generally leads to better outcomes. Many conditions can be effectively managed with appropriate treatment. Regular follow-up care is important for monitoring and treatment adjustment.
